\clearpage\fancyhead[L]{\small VESTA.}
\begin{center}
VESTA.
\end{center}

\begin{center}
\rule{0.6\linewidth}{0.4pt}\\
Monatliche Correspondenz. Band XV. S. 590–600. Juni 1807.\\
\rule{0.6\linewidth}{0.4pt}
\end{center}

--- --- ---

\par\indent Herr Dr. \textsc{Gauss} in Braunschweig machte die folgenden zwei Beobachtungen, die auf zahlreichen Vergleichen mit dem Stern von \textsc{Piazzi} beruhen, Rektaszension $180^\circ\,51'\,27''5$, nördliche Deklination $12^\circ\,35'\,14''0$.

\begin{center}\begingroup
\fontsize{10}{12}\selectfont
\setlength{\tabcolsep}{2.5pt}\renewcommand{\arraystretch}{1.15}
\begin{tabular}{c|c|c|c}
1807. & \shortstack{Mittl. Zeit in\\Braunschweig} & \shortstack{Scheinbare\\Rektaszension \VestaRATwoNinety{}} & \shortstack{Scheinb. nördliche\\Deklination \VestaDecTwoNinety{}} \\ \hline
April 6 & $8^{\text{u}}\,22^{\text{m}}\,45^{\text{s}}$ & $182^\circ\,21'\,54''6$ & $12^\circ\,28'\,0''9$ \\
8 & $8\ \ 28\ \ 49$ & $181\ \ 56\ \ 51.0$ & $12\ \ 35\ \ 45.8$
\end{tabular}\endgroup\end{center}

\par\indent Es war zu erwarten, dass der bewundernswerte Dr. \textsc{Gauss}, der sich um alle neu entdeckten kleinen Planeten so einzigartige Verdienste erworben hat, auch bei \emph{Vesta} unübertroffen bleiben würde. Doch diesmal hat sich dieser große Geometer selbst übertroffen. „Unseren vortrefflichen \textsc{Gauss} habe ich schon oft bewundert“, (schreibt uns Dr. \textsc{Olbers} am 20. April) „aber diesmal hat er mich doch in Erstaunen versetzt. Er erhielt meine Beobachtung der \emph{Vesta} vom 17. April am Abend des 20. April um $9\tfrac{1}{2}$ Uhr. Schon am 21. April um 5 Uhr nachmittags konnte er mir die berechneten elliptischen Elemente und den Vergleich mit allen Beobachtungen schicken. Für all dies-
\SourceLeaf{291}%
brauchte er nur zehn Stunden. Das belegt zugleich die große Flexibilität seiner neuen, so sehr verbesserten Methode\HistoricalNote{27} zur Bestimmung von Planetenbahnen.“

\par\indent Dr. \textsc{Gauss} selbst schreibt uns darüber Folgendes. „Die erste Bestimmung der Bahn nahm ich am 20. April vor, als ich die Beobachtung von Dr. \textsc{Olbers} vom 17. April erhalten hatte. Dadurch war der Zeitraum auf 19 Tage angewachsen. Die Ergebnisse übersandte ich Dr. \textsc{Olbers} am 21. April. Doch der Zeitraum war noch so kurz, fast alle Beobachtungen waren nur mit dem Kreismikrometer gemacht worden, und die geringe Neigung der Bahn dieses Planeten trägt außerdem dazu bei, die Ergebnisse vergleichsweise weniger zuverlässig zu machen. Deshalb wagte ich es noch nicht, sie schon als wirkliche Annäherung an die Wahrheit auszugeben, und hielt sie in der Absicht zurück, sie zunächst noch einmal anhand späterer Beobachtungen zu verbessern. Erst gestern (27. April) habe ich neuere Beobachtungen erhalten, mit denen ich gleich eine neue Berechnung anstellen wollte. Zu meiner großen Freude finde ich jedoch, dass meine ersten Ergebnisse noch so genau mit ihnen übereinstimmen, dass ich eigentlich noch gar nichts daran mit Sicherheit zu ändern wüsste. Ich schiebe diese neue Berechnung also auf, bis der Fehler eine deutlich merkliche Größe erreicht hat, und warte nun nicht länger damit, Ihnen die I. Elemente der \emph{Vesta} mitzuteilen.

\begin{center}\begingroup\fontsize{9}{11}\selectfont
\begin{tabular}{ll}
\hypertarget{vesta1807-first-elements}{}\shortstack[l]{Epoche der Länge, 29. März 1807\\$12^{\text{u}}$ M.\,Z. in Bremen} . & $193^\circ\,8'\,4''6$ \\
Aphel .\ .\ .\ .\ .\ .\ .\ .\ .\ .\ .\ .\ .\ .\ .\ . & $69^\circ\;7\;40.9$ \\
Aufsteigender Knoten .\ .\ .\ .\ .\ .\ .\ .\ .\ .\ .\ .\ . & $103\;8\;36.2$
\end{tabular}\endgroup\end{center}

\noindent Alle drei werden vom mittleren Äquinoktium aus gezählt, die beiden letzten gelten als tropisch ruhend.

\begin{center}\begingroup\fontsize{9}{11}\selectfont
\begin{tabular}{ll}
Mittlere tägliche Bewegung .\ .\ .\ .\ .\ .\ .\ .\ .\ .\ .\ . & $978''909$ \\
Exzentrizität .\ .\ .\ .\ .\ .\ .\ .\ .\ .\ .\ .\ .\ .\ .\ .\ .\ . & $0.097505$ \\
Logarithmus der großen Halbachse .\ .\ .\ .\ .\ .\ .\ .\ . & $0.3728428$ \\
Neigung der Bahn .\ .\ .\ .\ .\ .\ .\ .\ .\ .\ .\ .\ .\ .\ . & $7^\circ\;5'\;49''5$
\end{tabular}\endgroup\end{center}

\par\indent „Hier nun der Vergleich mit sämtlichen bekannt gewordenen Beobachtungen.

\begin{center}\begingroup
\fontsize{8}{10}\selectfont
\setlength{\tabcolsep}{2.5pt}\renewcommand{\arraystretch}{1.15}
\begin{tabular}{c|c|c|c|c|c|c}
1807 & \shortstack{Mittlere Zeit\\am\\Beobachtungsort} & \shortstack{Berechnete\\Rektaszension\\\VestaRATwoNinetyOne{}} & \shortstack{Unterschied.} & \shortstack{Berechnete nördl.\\Deklination\\\VestaDecTwoNinetyOne{}} & \shortstack{Unterschied.} & Beobachter. \\ \hline
März 29 & $8^{\text{u}}\,21^{\text{m}}\,20^{\text{s}}$ & $184^\circ\,8'\,54''4$ & $-1''4$ & $.\ .\ .\ .\ .$ & $.\ .\ .$ & Olbers \\
        & $10\ \ 31\ \ 16$ & $184\ \ 7\ \ 40.2$ & $-10.6$ & $11^\circ\,47'\,52''4$ & $+5''4$ & \\
30 & $8\ \ 34\ \ 53$ & $183\ \ 55\ \ \ 3.0$ & $+6.2$ & $.\ .\ .\ .\ .$ & $.\ .\ .$ & \\
   & $8\ \ 44\ \ \ 8$ & $183\ \ 54\ \ 57.7$ & $+11.9$ & $11\ \ 53\ \ 18.6$ & $+7.6$ & \\
   & $12\ \ 33\ \ 17$ & $183\ \ 52\ \ 44.2$ & $+3.4$ & $11\ \ 54\ \ 14.7$ & $-12.3$ & \\
April 1 & $9\ \ 50\ \ \ 0$ & $183\ \ 27\ \ \ 3.7$ & $+0.9$ & $12\ \ \ 4\ \ 40.7$ & $-11.3$ & \\
   & $12\ \ \ 4\ \ 35$ & $183\ \ 25\ \ 47.8$ & $+23.6$ & $12\ \ \ 5\ \ \ 9.6$ & $-22.1$ & Bessel \\
2 & $8\ \ 21\ \ \ 1$ & $183\ \ 14\ \ 23.8$ & $+6.0$ & $12\ \ \ 9\ \ 34.7$ & $-12.3$ & Olbers \\
3 & $8\ \ 16\ \ 49$ & $183\ \ \ 1\ \ \ 3.6$ & $+2.8$ & $12\ \ 14\ \ 33.6$ & $-6.4$ & \\
4 & $9\ \ \ 9\ \ 43$ & $182\ \ 47\ \ 20.9$ & $+0.1$ & $12\ \ 19\ \ 29.1$ & $-15.9$ & \\
5 & $11\ \ 17\ \ \ 3$ & $182\ \ 33\ \ 10.7$ & $-0.2$ & $12\ \ 24\ \ 21.6$ & $+2.5$ & Harding \\
6 & $8\ \ 22\ \ 45$ & $182\ \ 21\ \ 48.5$ & $-6.1$ & $12\ \ 28\ \ \ 6.6$ & $+5.7$ & Gauss \\
  & $11\ \ 12\ \ 15$ & $182\ \ 20\ \ 21.6$ & $-12.2$ & $12\ \ 28\ \ 34.6$ & $+45.2$ & Harding \\
8 & $8\ \ 21\ \ 37$ & $181\ \ 56\ \ 42.9$ & $-0.1$ & $12\ \ 35\ \ 50.1$ & $-8.0$ & Olbers \\
  & $8\ \ 28\ \ 49$ & $181\ \ 56\ \ 42.2$ & $-8.8$ & $12\ \ 35\ \ 50.1$ & $+4.3$ & Gauss \\
  & $11\ \ \ 2\ \ 43$ & $181\ \ 55\ \ 22.5$ & $+3.3$ & $12\ \ 36\ \ 13.6$ & $+1.4$ & Harding
\end{tabular}\endgroup\end{center}
\begin{center}\begingroup
\fontsize{8}{10}\selectfont
\setlength{\tabcolsep}{2.5pt}\renewcommand{\arraystretch}{1.15}
\begin{tabular}{c|c|c|c|c|c|c}
\SourceLeaf{292}%
1807 & \shortstack{Mittlere Zeit\\am\\Beobachtungsort} & \shortstack{Berechnete\\Rektaszension\\\VestaRATwoNinetyTwo{}} & \shortstack{Unterschied.} & \shortstack{Berechnete nördl.\\Deklination\\\VestaDecTwoNinetyTwo{}} & \shortstack{Unterschied.} & Beobachter. \\ \hline
April 9 & $8^{\text{u}}\,30^{\text{m}}\,46^{\text{s}}$ & $181^\circ\,44'\,25''6$ & $+8''6$ & $12^\circ\,39'\,19''8$ & $+5''7$ & Olbers \\
   & $10\ \ 57\ \ 58$ & $181\ \ 43\ \ 13.8$ & $+11.5$ & $12\ \ 39\ \ 39.6$ & $+3.5$ & Harding \\
12 & $8\ \ 27\ \ 20$ & $181\ \ \ 9\ \ 25.6$ & $-8.4$ & $12\ \ 48\ \ \ 7.4$ & $-8.7$ & Olbers \\
13 & $10\ \ 39\ \ 14$ & $180\ \ 57\ \ 28.3$ & $+11.3$ & $12\ \ 50\ \ 41.5$ & $-1.5$ & Bode \\
14 & $8\ \ 19\ \ 51$ & $180\ \ 47\ \ 36.3$ & $-21.7$ & $12\ \ 52\ \ 37.2$ & $-10.9$ & Olbers \\
   & $10\ \ 34\ \ 33$ & $180\ \ 46\ \ 39.2$ & $-0.7$ & $12\ \ 52\ \ 44.4$ & $+4.7$ & Harding \\
17 & $8\ \ 12\ \ 35$ & $180\ \ 17\ \ 21.6$ & $-5.2$ & $12\ \ 57\ \ 18.3$ & $+2.3$ & Olbers \\
23 & $8\ \ 30\ \ 52$ & $179\ \ 26\ \ 39.7$ & $-1.3$ & $12\ \ 59\ \ 20.4$ & $+0.4$ & \\
25 & $9\ \ 53\ \ 49$ & $179\ \ 13\ \ \ 0.6$ & $-18.4$ & $12\ \ 57\ \ 52.4$ & $-5.6$ & \\
   & $9\ \ 49\ \ 26$ & $179\ \ 12\ \ 57.1$ & $+8.1$ & $12\ \ 57\ \ 51.9$ & $-14.1$ &
\end{tabular}\endgroup\end{center}

\par\indent Diese Übereinstimmung ist so gut, wie man es bei diesen Beobachtungen nur erwarten kann. Da der Zeitraum nun 27 Tage beträgt, meint Dr. \textsc{Gauss}, dass diese Elemente wenigstens eine Vorstellung von den wichtigsten Abmessungen der Bahn geben. \emph{Vesta} würde demnach ihre Bahn viel schneller durchlaufen als die anderen kleinen Planeten. Ihre Bahnneigung wäre von allen die kleinste, ihre Exzentrizität mäßig. Letzteres und die große Helligkeit des Planeten geben Dr. \textsc{Gauss} sehr große Hoffnung, dass wir ältere Beobachtungen dieses neuen Planeten finden werden, vielleicht sogar bei \textsc{Flamsteed}. Sobald die Elemente etwas zuverlässiger bekannt sind, wird Dr. \textsc{Gauss} danach suchen.

\par\indent Am 18. Mai schreibt uns dieser unvergleichliche Astronom. „Meine ersten Elemente der \emph{Vesta}, die ich Ihnen vor einiger Zeit geschickt hatte und die auf Beobachtungen über 19 Tage beruhten, stimmten bis gegen Ende April durchweg sehr gut mit den Beobachtungen überein. Im Mai begann der Unterschied aber einigermaßen merklich zu werden und war am 6. Mai auf $+33''$ in der Rektaszension und $-41''$ in der Deklination angewachsen. Eine Verbesserung der Elemente wäre noch kaum nötig gewesen. Der Hauptgrund war, dass mir einige Sterne aus älteren Katalogen verdächtig geworden waren und ich ungeduldig wissen wollte, wie groß die noch nötigen Änderungen an den Elementen etwa sein könnten. Am 9. Mai, als ich die Beobachtung von Dr. \textsc{Olbers} vom $6^{\text{ten}}$ erhalten hatte, nahm ich eine zweite Berechnung der Bahn vor, deren Ergebnisse ich Ihnen hier mitzuteilen die Ehre habe.

\Needspace{16\baselineskip}
\begin{center}
II. \emph{Elemente der Vesta}
\end{center}

\begin{center}
Auf Beobachtungen über 38 Tage gegründet.
\end{center}

\begin{center}\begingroup\fontsize{9}{11}\selectfont
\begin{tabular}{ll}
\shortstack[l]{Epoche der Länge, 29. März 1807 $12^{\text{u}}$\\M.\,Z. in Bremen} .\ .\ .\ .\ .\ . & $192^\circ\,9'\,53''9$ \\
Aphel .\ .\ .\ .\ .\ .\ .\ .\ . & $69\;57\;52$ \\
Aufsteigender Knoten .\ .\ .\ .\ .\ . & $103\;18\;34$ \\
Mittlere tägliche tropische Bewegung  & $980''707$ \\
Exzentrizität .\ .\ .\ .\ .\ .\ .\ .\ . & $0.0872230$ \\
Logarithmus der mittleren Entfernung . & $0.3723521$ \\
Neigung der Bahn .\ .\ .\ .\ .\ .\ . & $7^\circ\,8'\,6''8$
\end{tabular}\endgroup\end{center}

\par\indent „Diese Elemente stimmen noch bis zum $11^{\text{ten}}$, also so weit, wie die mir vorliegenden Beobachtungen reichen, aufs Schönste mit ihnen überein. Ich habe bisher insgesamt 68 Beobachtungen verglichen. Die Übereinstimmung ist durchweg so gut, dass sich nirgends ein deutlicher Unterschied zwischen dem berechneten und dem wahren Ort angeben lässt.“

\par\indent „Meine Hoffnung, dass wir Beobachtungen der \emph{Vesta} in der \emph{Historia coelestis} finden werden, ist immer noch sehr groß. Bei den vermissten Sternen, die ich gleich am Anfang als verdächtig vermerkt hatte, ist diese Hoffnung aber größtenteils geschwunden. Nur über 91 \emph{Virginis} nach
\SourceLeaf{293}%
\textsc{Flamsteed} will ich noch kein abschließendes Urteil fällen. Die Länge würde sehr gut zum Ort der \emph{Vesta} passen, wenn die mittlere tägliche Bewegung nur etwa $6''$ größer wäre als nach den II. Elementen. Die Breite \emph{scheint} allerdings etwas zu groß. Bald werden wir darüber bestimmter urteilen können.“

--- --- ---
\hfill [\textsc{von Zach}.]
